\CWHeader{Лабораторная работа \textnumero 3}

\CWProblem{
Для реализации словаря из предыдущей лабораторной работы требуется провести исследование качества программы: оценить скорость выполнения и потребление оперативной памяти, а при обнаружении ошибок или явных недочётов исправить их.

В данной работе объектом исследования является словарь на основе красно-чёрного дерева, реализованный в лабораторной работе \textnumero 2.

Результатом работы должны являться:
\begin{enumerate}
\item дневник выполнения исследования;
\item выводы о найденных недочётах;
\item сравнение исправленной версии программы с исходной;
\item общие выводы о проделанной работе и полученном опыте.
\end{enumerate}

В качестве инструментов анализа допускается использовать не только \texttt{gprof} и \texttt{dmalloc}, но и их аналоги. В данной работе использовались доступные в macOS системные утилиты \texttt{sample}, \texttt{leaks} и \texttt{heap}, а также собственные измерительные программы для получения повторяемых временных характеристик.
}
\pagebreak
